\begin{table}[t]
\centering
\caption{Progressive component ablation of sparse-prompt CTV completion.}
\label{tab:ablation-summary}
\begin{threeparttable}
\resizebox{\linewidth}{!}{
\begin{tabular}{@{}clcccccc@{}}
\toprule
Step & Variant & Full-slice prompt & SDF support & Core/envelope & Supervised refine & DSC & uDSC \\
\midrule
1 & Linear mask interpolation & \checkmark & -- & -- & -- & $0.912 \pm 0.048$ & $0.876 \pm 0.057$ \\
2 & SDF propagation & \checkmark & \checkmark & -- & -- & $0.916 \pm 0.045$ & $0.871 \pm 0.064$ \\
3 & SDF core & \checkmark & \checkmark & \checkmark & -- & $0.920 \pm 0.045$ & $0.879 \pm 0.061$ \\
4 & Support-intersection pseudo-label & \checkmark & \checkmark & \checkmark & -- & $0.928 \pm 0.050$ & $0.897 \pm 0.062$ \\
5 & \textbf{Ours: pseudo-to-true refine network} & \checkmark & \checkmark & \checkmark & \checkmark & $\mathbf{0.934 \pm 0.045}$ & $\mathbf{0.906 \pm 0.054}$ \\
\bottomrule
\end{tabular}
}
\begin{tablenotes}[flushleft]\footnotesize
\item The table is ordered by the workflow rather than by unrelated parameter sweeps. Step 4 is the input pseudo-label used by the network. Step 5 adds supervised cohort-level pseudo-to-true correction without patient-specific training at test time.
\end{tablenotes}
\end{threeparttable}
\end{table}
